\section*{الفصل \ref{ch:b3:bacteriology} --- علم الجراثيم: النمو والفيزيولوجيا والوراثة}

\begin{solution}{exo:b3:bacteriology:1}
إيجابية الغرام: غشاء هيولي ملفوف بجدار \omterm{def:b3:bacteriology:envelope}{ببتيدوغليكان} ثخين
متعدد الطبقات تتخلله حموض التيكويك. وسلبية الغرام: طبقة
\omterm{def:b3:bacteriology:envelope}{ببتيدوغليكان} رقيقة في محيط حول البلازما، ثم \omterm{def:b3:bacteriology:envelope}{غشاء خارجي} وريقته
الخارجية \omterm{def:b3:bacteriology:envelope}{عديد سكريد شحمي}، مع بورينات. ومعقّد البنفسجي البلوري واليود
محبوس في الجدار الثخين المجفَّف بالكحول عند
إيجابية الغرام (فتكون بنفسجية) لكنه يُغسل من سلبية الغرام عبر جدارها
الرقيق وغشائها الخارجي المفكَّك، فتأخذ حينئذ الصبغة
المضادة الوردية.
\end{solution}

\begin{solution}{exo:b3:bacteriology:2}
زيادة $10^{6}$ ضعف $= 2^{19.9}$: أي نحو $20$ جيلًا؛ و$T = \qty{8}{h}/20
= \qty{24}{min}$؛ و$\mu = \ln 2/\qty{0.4}{h} = \qty{1.7}{h^{-1}}$.
\end{solution}

\begin{solution}{exo:b3:bacteriology:3}
\omterm{def:b3:bacteriology:hgt}{التحول} --- غريفيث (1928) وآفري وماكلاود وماكارتي
(1944): فالمكورات الرئوية الفوعة المقتولة بالحرارة تحوّل
الحية غير الفوعة، والعامل DNA. \omterm{def:b3:bacteriology:hgt}{والنقل بالعاثيات} --- زيندر وليدربرغ
(1952): فالعاثية P22 تحمل مورّثات \emph{Salmonella} من سلالة إلى
أخرى عبر مرشح يوقف الخلايا. \omterm{def:b3:bacteriology:hgt}{والاقتران} --- ليدربرغ
وتاتوم (1946): فسلالتان معوزتان من \emph{E.~coli} تعطيان
مؤتلفات مكتفية حين تُخلطان فقط، وهو ما يقتضي تماسًا خلويًا
(أنبوب دافيس على شكل U).
\end{solution}

\begin{solution}{exo:b3:bacteriology:4}
البنسلين: ناقلات الببتيد التي تشبّك \omterm{def:b3:bacteriology:envelope}{الببتيدوغليكان}، وهي ما
لا تملكه خلايا الحيوان. والستربتوميسين: الوحدة الريبوزومية الصغيرة،
وصورتها البكتيرية تختلف عن صورة حقيقيات النوى. والسيبروفلوكساسين: جيراز
DNA، وهو توبوإيزوميراز بكتيري غائب عن البشر. والريفامبيسين:
الوحدة $\beta$ من بوليميراز RNA البكتيري، وهي ما لا يشاركه فيه
إنزيم حقيقيات النوى.
\end{solution}

\begin{solution}{exo:b3:bacteriology:5}
$S^{*} = K_{S}D/(\mu_{\max} - D)$، و$X^{*} = Y(S_{0} - S^{*})$. فعند $D = 0.2$:
$S^{*} = \qty{0.0067}{g/L}$، و$X^{*} = \qty{0.80}{g/L}$. وعند $D = 0.6$:
$S^{*} = 0.06$، و$X^{*} = 0.78$. وعند $D = 0.79$: $S^{*} = 1.58$، و$X^{*} = 0.17$.
والغسل عند $D = \mu(S_{0}) = 0.8\times 2/2.02 = \qty{0.79}{h^{-1}}$.
\end{solution}

\begin{solution}{exo:b3:bacteriology:6}
$(N/V)_{\text{thr}} = kA_{\text{thr}}/p$. في عضو الضوء: $0.01\times
6\times 10^{12}/100 = 6\times 10^{8}$ في اللتر $= 6\times 10^{5}$ في
المليلتر. وفي ماء البحر: $10\times 6\times 10^{12}/100 = 6\times 10^{11}$ في اللتر
$= 6\times 10^{8}$ في المليلتر --- أي أعلى بألف مرة، ولا تُبلغ أبدًا في
البحر.
\end{solution}

\begin{solution}{exo:b3:bacteriology:7}
\qty{25}{mm}: حساسة، وMIC منخفض. و\qty{12}{mm}: حساسية منخفضة،
متوسطة --- وقد تخفق عند الجرعات العادية. و\qty{0}{mm}: مقاومة،
والنمو حتى القرص. والدواء ينتشر من القرص فيهبط
تركيزه مع المسافة هبوطًا قابلًا للتكرار؛ ويتوقف النمو
حيث يساوي التركيز MIC، فيرتبط نصف قطر المنطقة وMIC
بمنحني معايرة أُقيم بسلالات معلومة MIC.
\end{solution}

\begin{solution}{exo:b3:bacteriology:8}
\omterm{def:b3:genetic-engineering:tools}{البلازميد} الاقتراني يحمل إنتغرونًا تحوي مصفوفة خزائنه
عدة مورّثات مقاومة، إضافةً إلى عناصر قافزة تحمل مزيدًا؛ \omterm{def:b3:genetic-engineering:tools}{والبلازميد}
يعبر في تزاوج واحد ويعبّر عنها كلها. وهي تتجمع لأن
الانتقاء لأي من الأدوية يحفظ \omterm{def:b3:genetic-engineering:tools}{البلازميد} كله، فتدوم
المورّثات التي تركبه (وهو الانتقاء المشترك)؛ وتلتقط الإنتغرونات
الخزائن؛ وتقفز العناصر القافزة على البلازميدات؛ \omterm{def:b3:genetic-engineering:tools}{والبلازميد}، لا
المورّثة، هو وحدة النقل.
\end{solution}

\begin{solution}{exo:b3:bacteriology:9}
الدواء وحده: $10^{-7}\times 10^{9} = 100$ خلية مقاومة متوقعة. وبدواءين
مستقلين: $10^{-14}\times 10^{9} = 10^{-5}$. أما المعنِّدات فليست
طافرات: فنسبة صغيرة من الخلايا الساكنة تنجو من أي دواء و
تعيد النمو جمهرةً حساسة تمامًا حين يزول الدواء، فلا يفيد دواء
ثانٍ؛ ولا يفيد إلا علاج يطول بما يكفي للإمساك بها وهي
تستيقظ، أو دواء فعال على الخلايا الساكنة.
\end{solution}

\begin{solution}{exo:b3:bacteriology:10}
عند $X > X^{*}$: يتجاوز الاستهلاك $\mu X/Y$ الإمداد، فتهبط $S$ دون
$S^{*}$؛ فتصير $\mu(S) < D$ و$\mathrm{d}X/\mathrm{d}t < 0$، فتهبط $X$
عائدةً؛ ومع هبوط $X$ ترتفع $S$ من جديد. فيُصحَّح الانحرافان معًا: أي إن
الحالة المستقرة مستقرة. وأما نوعان: فيحتاج كل منهما إلى $\mu_{i}(S) = D$ ليدوم،
وهذا يثبّت $S_{i}^{*}$ الخاص به؛ ولا يستطيع المغذي أن يأخذ إلا
قيمة واحدة، فلا يدوم أكثر من نوع واحد. والذي له $S^{*}(D)$ الأدنى
يدفع $S$ دون ما يحتاجه الآخر، فيُغسل الآخر خارجًا؛ وأي النوعين
هو ذاك قد يتغير مع $D$ إذا تقاطع منحنيا مونو عندهما.
\end{solution}

\begin{solution}{exo:b3:bacteriology:11}
بالكثافة $\rho = N/V$، توجد الحالة المغلقة $A = p_{0}\rho/k$ ما دامت
تبقى دون $A_{\text{thr}}$، أي $\rho < kA_{\text{thr}}/p_{0}$؛
وتوجد الحالة المفتوحة $A = p_{1}\rho/k$ ما دامت تبقى فوقها، أي $\rho
> kA_{\text{thr}}/p_{1}$. ومن أجل $kA_{\text{thr}}/p_{1} < \rho <
kA_{\text{thr}}/p_{0}$ تتسق الحالتان مع نفسيهما: فالجمهرة التي
فتحت تبقي نفسها مفتوحة بناتجها الأعلى، والتي لم تفتح
تبقى مغلقة. والفائدة: الالتزام --- فمتى قررت مستعمرة أن تبني
غشاءً حيويًا أو تطلق ذيفانًا، لم ينقض هبوطٌ في الكثافة القرار،
ويكون المفتاح حادًا متزامنًا لا متدرجًا.
\end{solution}

\begin{solution}{exo:b3:bacteriology:12}
(أ) جرعات عالية لأشواط قصيرة: يبقى المستوى فوق MIC الطافر
ويقتل النمط البري وطافرات الخطوة الأولى سواء بسواء --- فمحبَّذة، وتحدّها
السمية. (ب) وجرعات منخفضة لأشواط طويلة: يقعد المستوى في
النافذة أيامًا فيربّي المقاومة --- فمرفوضة. (ج) والتوقف حين
تزول الأعراض: فالناجون أقل الخلايا حساسية والمستوى الهابط
يمر بالنافذة --- فمرفوض. (د) والتوليفة:
على الخلية أن تكون مقاومة للاثنين، وهو ما لا يحويه مليار خلية
--- فمحبَّذة.
\end{solution}

\begin{solution}{pb:b3:bacteriology:1}
\textbf{1.} $18$ جيلًا: أي $2^{18} = 2.6\times 10^{5}$ خلية،
و$2.6\times 10^{3}$ في المليلتر.
\textbf{2.} $2\times 10^{9}\times 100 = 2\times 10^{11}$ خلية $=
2^{37.5}$: أي بعد \qty{12.5}{h}؛ والكتلة \qty{0.2}{g}.
\textbf{3.} $6\times 10^{27}/10^{-12} = 6\times 10^{39} = 2^{132}$ خلية:
أي $132$ جيلًا، و\qty{44}{h}. والمغذيات والأكسجين والفضلات توقفه
في يوم؛ فالنمو الأسّي قصير دائمًا.
\textbf{4.} خلايا \omterm{def:b3:bacteriology:phases}{الطور الثابت} قد فككت ريبوزوماتها وكبتت
إنزيماتها الاستقلابية؛ فعليها أن تعيد بناءها قبل أن تنقسم.
أما الخلايا الأسّية فتصل مجهزة تمامًا.
\textbf{5.} $\mu = \ln 2/\qty{0.333}{h} = \qty{2.08}{h^{-1}}$؛ وعند $T =
\qty{60}{min}$، تكون \qty{0.69}{h^{-1}}؛ وبعد \qty{6}{h} يكون $2^{6} = 64$
بدل $2^{18}$: أي أقل بعامل $2^{12} \approx 4000$.
\textbf{6.} $200/(0.1\times 10^{-6}) = 2\times 10^{9}$ خلية حية في
المليلتر. والمجهر يعدّ أيضًا الخلايا الميتة والتي لن تنمو
على الطبق، والكتلة الواحدة تعطي مستعمرة واحدة.
\textbf{7.} $D = 0.5/1 = \qty{0.5}{h^{-1}}$؛ و$T = \ln 2/0.5 =
\qty{1.39}{h} = \qty{83}{min}$.
\textbf{8.} $S^{*} = 0.01\times 0.5/0.5 = \qty{0.01}{g/L}$؛ و$X^{*} =
0.5\times 1.99 \approx \qty{1}{g/L} = 10^{12}$ خلية في اللتر، أي $10^{9}$
في المليلتر.
\textbf{9.} $0.5\times 10^{12} = 5\times 10^{11}$ خلية في الساعة؛ و$720/1.39
\approx 520$ جيلًا في الشهر.
\textbf{10.} عند $D = 0.95$: $S^{*} = 0.01\times 0.95/0.05 = \qty{0.19}{g/L}$،
و$X^{*} = 0.5\times 1.81 = \qty{0.9}{g/L}$. وعند $F = \qty{1.1}{L/h}$، يكون $D
> \mu_{\max}$: أي غسل.
\textbf{11.} الطافر يبقي $S^{*} = 0.01\times 0.5/0.7 = \qty{0.0071}{g/L}$؛
وعند ذلك الغلوكوز تنمو الأصلية بمعدل $1.0\times 0.0071/0.0171 =
\qty{0.42}{h^{-1}} < D$ وتتناقص مثل $e^{-0.08t}$: فتُغسل خارجًا على مدى
أسابيع قليلة. ويستولي الطافر.
\textbf{12.} $10^{-9}\times 10^{12} = 1000$ طفرة نافعة في
الجيل: فالتكيف مستمر، بنسيلة أفضل جديدة
تجتاح كل بضع مئات من الأجيال --- فالمزرعة الكيميائية
آلة تطور.
\textbf{13.} $kA_{\text{thr}}/p = 0.05\times 6\times 10^{12}/500 =
6\times 10^{8}$ في اللتر $= 6\times 10^{5}$ خلية في المليلتر.
\textbf{14.} $10^{10}/6\times 10^{5} \approx 1.7\times 10^{4}$ ضعفًا.
\textbf{15.} $20\times 6\times 10^{12}/500 = 2.4\times 10^{11}$ في اللتر،
أي $2.4\times 10^{8}$ في المليلتر --- أي ستة رتب فوق
$10^{2}$ في المليلتر في البحر: فمظلمة.
\textbf{16.} عشرة أضعاف هي $3.3$ مضاعفة، أي \qty{6.6}{h}. وعند $10^{9}$ في المليلتر
تكون الجمهرة ما تزال فوق العتبة بمقدار $1700$ ضعف: فيبقى الضوء
مفتوحًا.
\textbf{17.} في البحر لا يرى الضوءَ أحد ولا يطعم مضيفًا: فالطافر
المظلم يوفر \qty{20}{\%} ويفوق المتوهجات نموًا. أما في العضو
فينتقي الحبار للضوء --- إذ يطرد الغشاشات المظلمة أو لا يطعمها
--- فتشتري الكلفة بيتًا.
\textbf{18.} احجب المستقبل أو اهدم \omterm{def:b3:bacteriology:quorum}{المحرّض الذاتي} (وهو إخماد
النصاب): فتعيش البكتيريا لكنها لا تنشر ذيفاناتها قط، ويصفّيها
الجهاز المناعي. ولأن البكتيريا الحساسة و``المقاومة'' تنمو
سواء --- إذ لا يُقتل شيء --- تكون أفضلية
الإفلات من الدواء صغيرة، وينبغي أن تنتشر المقاومة أبطأ
منها تجاه دواء قاتل.
\textbf{19.} للدواء A: $10^{-8}\times 10^{9} = 10$؛ وللدواء B: $100$؛ وللاثنين: $10^{-15}\times
10^{9} = 10^{-6}$.
\textbf{20.} $P_{0}(\text{A}) = e^{-10} = 5\times 10^{-5}$؛
و$P_{0}(\text{كليهما}) = e^{-10^{-6}} = 0.999999$.
\textbf{21.} $16 \to 8$: عمر نصف واحد، أي \qty{4}{h}؛ و$16 \to 1$: أربعة،
أي \qty{16}{h}؛ فتكون في النافذة من \qty{4}{h} إلى \qty{16}{h}، أي \qty{12}{h}
لكل جرعة.
\textbf{22.} بجرعة كل \qty{12}{h}: تكون في النافذة من الساعة $4$ إلى
الساعة $12$، أي ثلثي الزمن؛ وعلى مدى عشرة أيام يُنتقى لطافرات
الخطوة الأولى للدواء A --- وهي عشر في البداية --- ويخفق A وحده.
\textbf{23.} لإبقائه فوق \qty{8}{mg/L} طوال الوقت: جرعة كل عمر نصف
(\qty{4}{h})، أو جرعة تبلغ ذروة \qty{64}{mg/L} كل \qty{12}{h}
($64 \to 8$ في ثلاثة أعمار نصف)، أو تسريب مستمر. والكلفة:
سمية الذرى العالية، أو ست جرعات في اليوم لا يلتزم بها
المرضى.
\textbf{24.} $10^{-4}\times 10^{9} = 10^{5}$ معنِّدة تنجو مهما كانت
الأدوية؛ فإذا توقف العلاج استيقظت وأعادت نمو جمهرة حساسة
تمامًا --- أي انتكاسة. فيجب أن يدوم العلاج بما يكفي للإمساك
بها وهي تغادر السكون، ولهذا يستغرق السل ستة أشهر.
\textbf{25.} $10^{9}$ خلية في المليلتر في المزرعة الكيميائية؛ ويضيء الضوء عند
$6\times 10^{5}$ خلية في المليلتر؛ و$P_{0} = 0.999999$ أن شوط الدوائين
لا يلقى خلية مقاومة.
\end{solution}
